\clearpage
\homeworktitle{Решения и комментарии}{Версия для учителя}

\begin{teacheranswer}[Задание 1]
  Дробь равна нулю, когда её числитель равен нулю, а знаменатель отличен
  от нуля. Из числителя получаем \(x=2\) и \(x=-3\). Но при \(x=2\)
  знаменатель равен нулю. Поэтому
  \[
    \boxed{x=-3}.
  \]
\end{teacheranswer}

\begin{criterionbox}
  Проверяется понимание того, что ноль числителя ещё не обязательно является
  корнем исходной дроби.
\end{criterionbox}

\begin{teacheranswer}[Задание 2]
  Раскрываем модуль:
  \[
    2|x|-x=
    \begin{cases}
      x, & x\ge 0,\\
      -3x, & x<0.
    \end{cases}
  \]
  При \(x\ge0\) получаем \(x=b\), поэтому требуется \(b\ge0\). При
  \(x<0\) получаем \(x=-\dfrac b3\), поэтому требуется \(b>0\).

  Следовательно,
  \[
    \begin{array}{c|c}
      b<0 & 0\text{ корней},\\
      b=0 & 1\text{ корень},\\
      b>0 & 2\text{ корня}.
    \end{array}
  \]
\end{teacheranswer}

\begin{criterionbox}
  Проверяется умение одновременно учитывать формулу ветви и условие,
  на котором эта ветвь рассматривается.
\end{criterionbox}

\clearpage
\section*{Решение задания 3}

\begin{teacheranswer}[Шаг 1. Исследуем числитель]
  Уравнение числителя имеет вид
  \[
    2|x|-x=a+1.
  \]
  При \(x\ge0\):
  \[
    \begin{cases}
      2x-x=a+1,\\
      x\ge0
    \end{cases}
    \quad\Longleftrightarrow\quad
    \begin{cases}
      x_1=a+1,\\
      a\ge-1.
    \end{cases}
  \]
  При \(x<0\):
  \[
    \begin{cases}
      -2x-x=a+1,\\
      x<0
    \end{cases}
    \quad\Longleftrightarrow\quad
    \begin{cases}
      x_2=-\dfrac{a+1}{3},\\
      a>-1.
    \end{cases}
  \]
  Поэтому числитель имеет два различных нуля тогда и только тогда, когда
  \[
    a>-1.
  \]
\end{teacheranswer}

\begin{teacheranswer}[Шаг 2. Проверяем знаменатель]
  Для \(x_1=a+1\):
  \[
    \begin{aligned}
      x_1^2-x_1-a
      &=(a+1)^2-(a+1)-a\\
      &=a^2.
    \end{aligned}
  \]
  Первый ноль числителя обращает знаменатель в ноль при \(a=0\).

  Для \(x_2=-\dfrac{a+1}{3}\):
  \[
    \begin{aligned}
      x_2^2-x_2-a
      &=\frac{(a+1)^2}{9}+\frac{a+1}{3}-a\\
      &=\frac{(a-2)^2}{9}.
    \end{aligned}
  \]
  Второй ноль числителя обращает знаменатель в ноль при \(a=2\).

  Значит,
  \[
    \boxed{a\in(-1,+\infty)\setminus\{0,2\}}.
  \]
\end{teacheranswer}

\clearpage
\section*{Решение задания 4}

\begin{teacheranswer}[Шаг 1. Исследуем числитель]
  Из уравнения
  \[
    3|x|-x=a+2
  \]
  получаем:
  \[
    x_1=\frac{a+2}{2}\quad(x\ge0),
    \qquad
    x_2=-\frac{a+2}{4}\quad(x<0).
  \]
  Оба нуля существуют и различны при
  \[
    a>-2.
  \]
\end{teacheranswer}

\begin{teacheranswer}[Шаг 2. Проверяем знаменатель]
  Для первого нуля:
  \[
    \begin{aligned}
      x_1^2-x_1-a
      &=\frac{(a+2)^2}{4}-\frac{a+2}{2}-a\\
      &=\frac{a(a-2)}4.
    \end{aligned}
  \]
  Поэтому исключаются \(a=0\) и \(a=2\).

  Для второго нуля:
  \[
    \begin{aligned}
      x_2^2-x_2-a
      &=\frac{(a+2)^2}{16}+\frac{a+2}{4}-a\\
      &=\frac{(a-2)(a-6)}{16}.
    \end{aligned}
  \]
  Поэтому исключаются \(a=2\) и \(a=6\).

  Окончательно
  \[
    \boxed{a\in(-2,+\infty)\setminus\{0,2,6\}}.
  \]
\end{teacheranswer}

\begin{criterionbox}
  Здесь важно заметить, что значение \(a=2\) возникает при проверке обоих
  нулей, но в окончательном ответе исключается только один раз.
\end{criterionbox}

\clearpage
\section*{Решения заданий 5 и 6}

\begin{teacheranswer}[Задание 5. Анализ ошибки]
  Условие \(a>-2\) необходимо и достаточно только для того, чтобы
  \emph{числитель} имел два различных нуля. Для исходной дроби нужно ещё
  проверить, не обращают ли эти значения знаменатель в ноль.

  Из разбора основной задачи известно, что при
  \[
    a\in\{-1,0,3,8\}
  \]
  один из нулей числителя не является корнем исходного уравнения. Поэтому
  правильный ответ:
  \[
    \boxed{a\in(-2,+\infty)\setminus\{-1,0,3,8\}}.
  \]
\end{teacheranswer}

\begin{teacheranswer}[Задание 6*. Исследуем числитель]
  Рассмотрим два случая.

  При \(x\ge a\):
  \[
    \begin{cases}
      2(x-a)-x-2=0,\\
      x\ge a
    \end{cases}
    \quad\Longleftrightarrow\quad
    \begin{cases}
      x_1=2a+2,\\
      a\ge-2.
    \end{cases}
  \]

  При \(x<a\):
  \[
    \begin{cases}
      2(a-x)-x-2=0,\\
      x<a
    \end{cases}
    \quad\Longleftrightarrow\quad
    \begin{cases}
      x_2=\dfrac{2a-2}{3},\\
      a>-2.
    \end{cases}
  \]

  Следовательно, числитель имеет два различных нуля при \(a>-2\).
  Действительно,
  \[
    x_1-x_2=\frac{4(a+2)}3>0.
  \]
\end{teacheranswer}

\begin{teacheranswer}[Задание 6*. Проверяем знаменатель]
  Знаменатель \(x^2-1\) обращается в ноль при \(x=\pm1\).

  Для \(x_1=2a+2\):
  \[
    2a+2=-1\Rightarrow a=-\frac32,
    \qquad
    2a+2=1\Rightarrow a=-\frac12.
  \]

  Для \(x_2=\dfrac{2a-2}{3}\):
  \[
    \frac{2a-2}{3}=-1\Rightarrow a=-\frac12,
    \qquad
    \frac{2a-2}{3}=1\Rightarrow a=\frac52.
  \]

  Все найденные значения лежат в промежутке \((-2,+\infty)\), поэтому их
  нужно исключить:
  \[
    \boxed{
      a\in(-2,+\infty)\setminus
      \left\{-\frac32,-\frac12,\frac52\right\}
    }.
  \]
\end{teacheranswer}

\begin{criterionbox}[Методический смысл задания со звёздочкой]
  Параметр находится внутри модуля, поэтому изменяется не уровень
  горизонтальной прямой, а положение точки излома графика. Задание проверяет
  перенос метода, а не воспроизведение готового алгоритма.
\end{criterionbox}

\begin{resultbox}[Краткие ответы]
  \begin{enumerate}[itemsep=0.28em,topsep=0.25em]
    \item \(x=-3\).
    \item \(b<0\): 0; \(b=0\): 1; \(b>0\): 2.
    \item \(a\in(-1,+\infty)\setminus\{0,2\}\).
    \item \(a\in(-2,+\infty)\setminus\{0,2,6\}\).
    \item \(a\in(-2,+\infty)\setminus\{-1,0,3,8\}\).
    \item \(a\in(-2,+\infty)\setminus
      \left\{-\dfrac32,-\dfrac12,\dfrac52\right\}\).
  \end{enumerate}
\end{resultbox}
